\section{Introduction}
\label{sec:introduction}

Standard New Keynesian (NK) models used by the Federal Reserve, the European
Central Bank, and the International Monetary Fund share a foundational
architecture: competitive labor markets clear wages at the marginal product of
labor, and firms source inputs from globally integrated, cost-minimizing supply
chains. Two decades of empirical accumulation and three years of post-pandemic
macroeconomic disruption have rendered both assumptions untenable as
approximations to the modern economy.

\paragraph{The Monopsony Fact.}
\citet{azar2022labor}, \citet{rinz2022labor}, and \citet{benmelech2022}
document that the median U.S.\ worker faces a local labor market
Herfindahl-Hirschman Index (HHI) exceeding 4{,}000---the Department of
Justice's threshold for ``highly concentrated.'' Wage markdowns (the gap
between the marginal product of labor and the observed wage) average 15--30\%
across sectors, with substantial cross-sectional variation positively
correlated with HHI \citep{azar2022labor}. Critically, this gap is
\emph{countercyclical}: firms compress wages more aggressively in downturns,
when workers' outside options deteriorate and firm-level labor supply
elasticities fall \citep{dube2020monopsony, manning2021monopsony}. This
decouples unemployment from inflation in ways standard models cannot reproduce,
and it implies that the aggregate Phillips curve is \emph{endogenously
flatter} than competition-based calibrations suggest.

\paragraph{The Fragmentation Fact.}
The IMF's Geoeconomic Fragmentation series \citep{aiyar2023} documents a
measurable bifurcation of global trade into politically aligned blocs,
accelerating sharply post-2022. Using bilateral trade flow data, \citet{gopinath2024}
show that trade within Western-aligned and China-aligned blocs has remained
resilient, while cross-bloc trade has been rerouted through costly
intermediary countries or reshored domestically. This structural shift means
supply shocks that were previously absorbed by flexible global substitution
are now amplified domestically and propagate through concentrated upstream
networks with elevated Domar weights \citep{baqaee2019}.

\paragraph{The Central Gap.}
No existing general equilibrium model integrates both of these facts in a
dynamic, quantitative framework. HANK-IO models \citep{auclert2021,
baqaee2019} handle heterogeneous agents and input-output networks but assume
competitive labor markets and integrated trade. Monopsony has been grafted
onto static or representative-agent NK frameworks
\citep{berger2022monopsony, autor2020fall}, but without the heterogeneous
agent redistribution channel or the multi-region network structure. This paper
fills that gap.

\subsection{What We Do}

We build a model with three novel, interacting components:

\begin{enumerate}
\item \textbf{Structural monopsony with a countercyclical markdown.}
  Firms face upward-sloping firm-level labor supply curves, parameterized by
  a firm-level elasticity $\varepsilon$ that falls when unemployment rises.
  The resulting wage markdown $\mu^w = \varepsilon/(\varepsilon+1)$ deepens in
  recessions, generating endogenous, state-dependent Phillips curve flattening.

\item \textbf{Multi-region, multi-bloc IO network.}
  The world is partitioned into geopolitical blocs. Trade costs between blocs
  are higher and subject to a time-varying \emph{fragmentation premium} $\phi_t$
  that follows an AR(1) process with high persistence. Firms source intermediates
  using a nested CES structure over domestic and foreign varieties; the
  Armington elasticity $\eta$ governs substitution across regions. We represent
  the resulting multi-sector production structure as a weighted directed graph
  and compute Domar weights and Leontief inverses analytically.

\item \textbf{Heterogeneous households with incomplete markets (HANK).}
  A continuum of households differs in asset positions and labor types. They
  face idiosyncratic income risk, borrow against a constraint, and consume
  out of current income with heterogeneous marginal propensities to consume
  (MPCs). The distribution of assets and labor types is a state variable.
\end{enumerate}

We solve the model using the \emph{sequence-space Jacobian} (SSJ) method of
\citet{auclert2021ssj}, extended with two novel Jacobian blocks:
(i) the \emph{monopsony block}, which adds the countercyclical markdown
response to the standard wage-employment Jacobian; and (ii) the
\emph{fragmentation block}, which maps shocks to the trade cost matrix
$\boldsymbol{\tau}(\phi_t)$ into shocks to the full $(S \times R)$-dimensional
Leontief inverse. Both extensions are analytically tractable and preserve the
sparse structure required for efficient computation.

\subsection{Main Results}

\paragraph{Proposition 1 (Phillips Curve Flattening).}
The monopsony-adjusted NKPC slope satisfies $\kappa^{mono} < \kappa^{comp}$,
with the gap proportional to the countercyclical markdown intensity $\gamma$:
\begin{equation*}
  \kappa^{mono} = \frac{\kappa^{comp}}{1 + \dfrac{\gamma}{\bar\varepsilon(\bar\varepsilon+1)} \cdot \dfrac{1}{\alpha_n}}.
\end{equation*}
In our baseline calibration ($\bar\varepsilon = 4$, $\gamma = 0.8$),
$\kappa^{mono}/\kappa^{comp} = 0.60$, implying a sacrifice ratio 67\% above the
competitive benchmark. This result holds without any additional price or wage
rigidity beyond the standard Calvo assumption.

\paragraph{Proposition 2 (Fragmentation Amplification).}
The semi-elasticity of aggregate TFP with respect to the fragmentation premium
is:
\begin{equation*}
  \frac{d\log\mathrm{TFP}}{d\phi} = -\sum_{\substack{(s,r),(s',r')\\b(r)\neq b(r')}}
    \lambda_{sr} \,\Omega_{(s,r)(s',r')} \,\frac{1}{\tau^{rr'}},
\end{equation*}
a closed-form, network-weighted expression. This quantity is strictly negative
and increasing in magnitude in the centrality of cross-bloc linkages---so
economies with higher upstream integration across blocs suffer larger TFP losses
from fragmentation. In our calibrated 10-sector, 5-region model, a 15pp
increase in $\phi$ reduces aggregate TFP by 1.2\% and raises inflation by
0.8pp, producing a stagflationary impulse that standard Taylor rules cannot
efficiently address.

\paragraph{Proposition 3 (Optimal Augmented Taylor Rule).}
The welfare-optimal monetary rule is:
\begin{equation*}
  i_t = r^* + \pi^* + \phi_\pi^*(\pi_t - \pi^*) + \phi_y^*\,\tilde{y}_t
        + \underbrace{\phi_\mu^*\,\hat\mu_t^w}_{\text{novel: }>0}
        + \underbrace{\phi_\phi^*\,\hat\phi_t}_{\text{novel: }<0} + \nu_t.
\end{equation*}
The sign of $\phi_\mu^* > 0$: when monopsony deepens in recession, compress
more of the income distribution below steady-state consumption, raising the
aggregate MPC and amplifying the recessionary drag on demand. The central bank
should ease more aggressively to offset this indirect channel. The sign of
$\phi_\phi^* < 0$: when fragmentation spikes, it is a cost-push shock---standard
Taylor rules tighten in response, amplifying the output loss without effectively
reducing supply-side inflation. The augmented rule correctly identifies this
as a non-demand shock and responds more mildly, achieving strictly lower
welfare loss.

\subsection{Relation to the Literature}

This paper contributes to three strands of literature.

\paragraph{Heterogeneous agent models.}
\citet{auclert2019} establishes the importance of redistribution channels for
monetary policy in HANK. \citet{kaplan2018} shows the dominance of indirect
effects. We add a new indirect channel: monopsony deepening in recessions
redistributes income toward firms and away from workers, elevating the aggregate
MPC and amplifying the recession for a given rate increase.

\paragraph{Input-output amplification.}
\citet{baqaee2019} and \citet{acemoglu2012} establish that sectoral shocks
propagate through networks weighted by Domar weights. We extend this to a
multi-bloc, multi-region dynamic setting where the Leontief inverse itself
varies with geopolitical conditions. \citet{caliendo2015} and
\citet{antras2020global} analyze trade in value chains; our contribution is
introducing a dynamic, monetary-policy-relevant version with time-varying
bloc structure.

\paragraph{Monopsony and the Phillips curve.}
\citet{berger2022monopsony} introduce monopsony into a static NK model.
\citet{autor2020fall} document the long-run labor market concentration trend.
\citet{manning2021monopsony} provides the key empirical estimates we use.
We are the first to integrate countercyclical monopsony into a HANK model
and derive the resulting implications for monetary policy transmission.

\subsection{Organization}

The paper proceeds as follows. Section \ref{sec:model} presents the model.
Section \ref{sec:equilibrium} defines equilibrium. Section \ref{sec:solution}
describes the solution method. Section \ref{sec:results} states and proves
the three main propositions. Section \ref{sec:quantitative} presents the
quantitative analysis. Section \ref{sec:policy} discusses policy implications.
Section \ref{sec:conclusion} concludes.
